\chapter{Discussion and Limitations}

\section{Main Findings}

\section{Sim-to-Real Gap}
7 — Discussion and Limitations
Proporrei tre sezioni, senza riprendere tutte le tabelle o le singole ablation.
7.1 Interpretation of the Main Findings
Due sottosezioni, dedicate ai due rami.
7.1.1 Lightweight Scene Understanding
- Interpretare il compromesso accuratezza–efficienza: SegFormer-B0 QAT privilegia l’accuratezza, mentre DeepLabV3+ MobileNetV2 QAT offre il minor consumo misurato.
- Spiegare perché parametri, memoria, latenza ed energia non sono indicatori intercambiabili.
- Discutere la quantizzazione come parte della configurazione di esecuzione: i risultati dipendono anche da backend, precisione mista e, per QAT, fine-tuning e distillazione.
- Ricavare dalle ablation una considerazione generale: augmentation e bilanciamento devono essere adatti ai dati, non aggiunti presumendo un miglioramento.
7.1.2 Synthetic Damage Inspection
- Discutere il diverso compromesso di PatchCore Large, SuperSimpleNet e TinyGLASS tra qualità e risorse.
- Evidenziare che detection e localization richiedono valutazioni distinte: riconoscere un’immagine danneggiata non significa delimitare accuratamente il foro.
- Interpretare il ruolo di crop, risoluzione, feature grid e aggregazione degli score.
- Mettere in evidenza la difficoltà dei piccoli difetti e la necessità di adattare i metodi al problema, anziché trasferire automaticamente configurazioni da altri benchmark.
Qui interessa spiegare che cosa abbiamo imparato dalle scelte, non raccontare nuovamente ogni esperimento.
7.2 Generalization and the Role of Data
Questa sezione collegherebbe i risultati alla diversa natura dei dati nei due rami.
- MarsSeg: l’accuratezza nella missione di training non garantisce trasferimento all’altra missione. SegFormer mostra il miglior trasferimento tra le configurazioni valutate, senza dimostrare una superiorità universale dei Transformer.
- Dataset sintetico: generazione controllata, maschere automatiche e coppie clean/damaged hanno reso possibile studiare il problema senza annotare un corpus reale.
- Domain randomization: introduce variabilità nel benchmark, ma i risultati disponibili non ne dimostrano da soli l’efficacia nel trasferimento al reale.
- Conoscenza del difetto: il training utilizza render normali, ma le perturbazioni hole-like di TinyGLASS incorporano un’ipotesi sull’aspetto del danno. Non abbiamo quindi dimostrato detection di qualunque anomalia sconosciuta.
- Distinguere chiaramente il trasferimento tra immagini reali delle missioni, effettivamente misurato, dal sim-to-real dell’inspection, ancora non verificato.
7.3 Methodological and Operational Limitations
Raccoglierei qui i limiti che circoscrivono le conclusioni:
- Confrontiamo configurazioni complete, con differenze di pretraining, risoluzione e preprocessing: non sempre possiamo isolare l’effetto dell’architettura.
- Le valutazioni presentate non quantificano la variabilità tra seed; alcuni sottogruppi hanno pochi campioni.
- L’inspection riguarda un asset di Curiosity, perforazioni visibili e, nel confronto principale, una posa controllata.
- Le metriche aggregate non stabiliscono una soglia operativa né un numero accettabile di falsi allarmi.
- I benchmark su T4 non dimostrano prestazioni su hardware di volo o consumi dell’intero sistema.
- I due rami non sono stati integrati; non abbiamo valutato navigazione autonoma, diagnosi meccanica o progressione del danno.
Per ogni limite, direi brevemente che cosa impedisce di concludere, senza anticipare la soluzione futura.
8 — Conclusions and Future Work
Qui manterrei due sezioni.
8.1 Conclusions
Un bilancio breve, organizzato attorno ai tre contributi fondamentali:
1. Scene understanding: confronto di modelli e strategie di quantizzazione, individuando configurazioni con diversi compromessi accuratezza–risorse.
2. Synthetic dataset generation: realizzazione della pipeline e del dataset annotato da 10.000 immagini per lo studio delle perforazioni.
3. Damage inspection: valutazione e adattamento di quattro famiglie di anomaly detector, distinguendo detection, localization ed efficienza.
Chiuderei spiegando che la tesi identifica soluzioni promettenti nelle condizioni studiate, senza presentarle come un sistema già validato su rover. Bastano pochi risultati rappresentativi: niente nuove spiegazioni delle ablation.
8.2 Future Work
Quattro direzioni, ordinate per priorità:
8.2.1 Real-Image Validation and Domain Adaptation
- Costruire un insieme annotato di immagini reali delle ruote e verificare il trasferimento dei detector.
- Studiare adattamento al dominio reale e nuove condizioni di acquisizione, anche per la segmentazione.
8.2.2 Broader Damage Coverage and Small-Defect Localization
- Estendere il dataset a più perforazioni, occlusioni, altri danni e asset.
- Investigare elaborazione locale o multiscala per conservare piccoli difetti senza aumentare indiscriminatamente il costo.
8.2.3 Integration of the Framework Branches
- Collegare la segmentazione a una localizzazione delle singole ruote e alla selezione automatica delle regioni da ispezionare.
- Valutare gli errori a cascata e il costo della pipeline completa.
8.2.4 Deployment and Operational Validation
- Verificare i modelli su hardware edge rappresentativo, includendo preprocessing e consumo complessivo.
- Esplorare quantizzazione dei detector, ripetizioni sperimentali e soglie compatibili con requisiti di falsi allarmi.
In sintesi: nel capitolo 7 spieghiamo perché un limite conta; nel capitolo 8 proponiamo come affrontarlo. Questa separazione evita che Discussion, Conclusions e Future Work diventino tre versioni dello stesso riepilogo.